\chapter{Carbonilas: adições nucleofílicas, acetais e proteção}\label{ch:b1:acetals-protection}

Dissolva glicose em água e quase nada dela permanece como o aldeído de
cadeia aberta desenhado nos livros: a molécula se dobra, seu próprio grupo
hidroxila se adiciona ao seu aldeído, e um anel de seis átomos se fecha. A
mesma adição de um álcool a um grupo carbonila, repetida mais uma vez com
um \omterm{def:b1:elementary-steps:catalyst}{catalisador} ácido, dá os \omterm{def:b1:acetals-protection:hemiacetal}{acetais}, compostos estáveis frente às bases e
aos \omterm{def:b1:electronic-effects:nucleophile}{nucleófilos} mais reativos. Essa estabilidade é uma ferramenta: um
aldeído que seria destruído por um \omterm{def:b1:organomagnesium:organometallic}{reagente de Grignard} pode ser
escondido como \omterm{def:b1:acetals-protection:hemiacetal}{acetal}, atravessar a reação e ser recuperado no fim. Este
capítulo estuda o grupo carbonila como \omterm{def:b1:electronic-effects:nucleophile}{eletrófilo}, as adições de água e
de álcoois, e a estratégia de \omterm{def:b1:acetals-protection:protecting-group}{proteção}.

\begin{recall}
As adições de Grignard aos compostos carbonílicos (\cref{ch:b1:organomagnesium});
o \omterm{def:b1:extent-q-and-k:quotient}{quociente de reação} $Q$ e a evolução de um sistema enquanto $Q \neq K$
(\cref{ch:b1:extent-q-and-k}); \omterm{def:b1:electronic-effects:curly-arrow}{setas curvas}, \omterm{def:b1:electronic-effects:nucleophile}{nucleófilos} e
\omterm{def:b1:electronic-effects:nucleophile}{eletrófilos} (\cref{ch:b1:electronic-effects}). O Livro 1 (ano 12)
apresentou a ideia de proteger um grupo durante uma síntese.
\end{recall}

\section{O grupo carbonila como eletrófilo}

A ligação C=O é polarizada no sentido do oxigênio (\cref{ch:b1:periodicity}),
e uma \omterm{def:b1:lewis-resonance-vsepr:resonance}{estrutura de ressonância} \ce{C+-O-} põe uma carga positiva inteira no
carbono: o carbono é eletrofílico, atacado pelos \omterm{def:b1:electronic-effects:nucleophile}{nucleófilos}
perpendicularmente ao plano do grupo. Os \omterm{def:b1:electronic-effects:nucleophile}{nucleófilos} fortes (\omterm{def:b1:organomagnesium:organometallic}{reagentes de
Grignard}, hidreto, cianeto) se adicionam diretamente. Os \omterm{def:b1:electronic-effects:nucleophile}{nucleófilos}
fracos --- água, álcoois --- precisam de ajuda.

\begin{definition}[Ativação eletrofílica]\label{def:b1:acetals-protection:activation}
A \emph{ativação eletrofílica}\index{ativacao eletrofilica@ativação eletrofílica} de um
grupo carbonila é sua protonação (ou sua coordenação a um \omterm{def:b1:lewis-resonance-vsepr:lewis-acid}{ácido de Lewis})
no oxigênio. O cátion \ce{C=OH+}, cuja carga positiva é compartilhada pelo
carbono por meio da \omterm{def:b1:lewis-resonance-vsepr:resonance}{estrutura de ressonância} \ce{C+-OH}, é atacado até
pelos \omterm{def:b1:electronic-effects:nucleophile}{nucleófilos} fracos.
\end{definition}

\begin{omfigure}
\setchemfig{atom sep=2.1em}
\schemestart
\chemfig{C(-[:150]R)(-[:210]H)=O}
\arrow{->[\ce{H+}]}
\chemfig{C(-[:150]R)(-[:210]H)=O^{+}-H}
\arrow{<->}
\chemfig{C^{+}(-[:150]R)(-[:210]H)-O-H}
\schemestop

{\small Ativação de um aldeído por um ácido: a protonação do oxigênio
produz um cátion cuja \omterm{def:b1:lewis-resonance-vsepr:resonance}{estrutura de ressonância} leva a carga no
carbono.}
\end{omfigure}

\section{Hidratos e hemiacetais}

A água se adiciona de modo reversível aos aldeídos e às cetonas, \ce{R2C=O + H2O <=>
R2C(OH)2}; o hidrato geralmente é minoritário para as cetonas e pode ser
majoritário para o metanal. Um álcool se adiciona da mesma maneira.

\begin{definition}[Hemiacetal e acetal]\label{def:b1:acetals-protection:hemiacetal}
Um \emph{hemiacetal}\index{hemiacetal} é um composto em que um mesmo
carbono leva um grupo OH e um grupo OR, formado pela adição de um álcool
a um grupo carbonila. Um \emph{acetal}\index{acetal} é um composto em que
um mesmo carbono leva dois grupos OR, formado a partir de um hemiacetal e
de uma segunda molécula de álcool, com perda de água.
\end{definition}

Os \omterm{def:b1:acetals-protection:hemiacetal}{hemiacetais} de cadeia aberta geralmente são minoritários em seu
equilíbrio com o aldeído e o álcool; quando o OH e a C=O pertencem à
mesma molécula e um anel de cinco ou seis átomos pode se fechar, o
\omterm{def:b1:acetals-protection:hemiacetal}{hemiacetal} cíclico predomina. É o caso da glicose, e do mais simples
5-hidroxipentanal.

\begin{omfigure}
\setchemfig{atom sep=2em}
\schemestart
\chemfig{HO-[:30]-[:-30]-[:30]-[:-30]-[:30]C(=[:90]O)-[:-30]H}
\arrow{<=>}
\chemfig{*6(O-(-[:-30]OH)-----)}
\schemestop

\medskip

{\small O 5-hidroxipentanal e seu \omterm{def:b1:acetals-protection:hemiacetal}{hemiacetal} cíclico, um anel de seis
átomos que contém o oxigênio. O OH do carbono 5 se adicionou ao aldeído
do carbono 1; a forma cíclica predomina, como para a glicose.}
\end{omfigure}

\section{Acetais}

\begin{proposition}[Mecanismo da formação dos acetais]\label{prop:b1:acetals-protection:acetal-mechanism}
Em meio ácido, um aldeído ou uma cetona e um álcool dão o \omterm{def:b1:acetals-protection:hemiacetal}{hemiacetal} e
depois o \omterm{def:b1:acetals-protection:hemiacetal}{acetal}, pelas etapas: protonação da C=O; adição do álcool;
perda de um próton (\omterm{def:b1:acetals-protection:hemiacetal}{hemiacetal}); protonação do OH; perda de água, que dá
um \omterm{def:b1:electronic-effects:intermediates}{carbocátion} estabilizado pelo grupo OR restante (um íon
oxocarbênio); adição de um segundo álcool; perda de um próton. Todas as
etapas são reversíveis.
\end{proposition}

\begin{proof}
Cada etapa é uma transferência de próton ácido--base ou a adição (ou a
perda) de um \omterm{def:b1:electronic-effects:nucleophile}{nucleófilo} a um (ou de um) carbono eletrofílico, como
descrito no \cref{ch:b1:electronic-effects}. A perda de água é possível
porque o cátion formado, $\mathrm{R_2C^+{-}OR}$, tem uma \omterm{def:b1:lewis-resonance-vsepr:resonance}{estrutura de ressonância} $\mathrm{R_2C{=}O^+R}$
em que todo átomo tem um octeto (\cref{prop:b1:electronic-effects:carbocation-order}).
O ácido é consumido nas protonações e regenerado nas
desprotonações: é um \omterm{def:b1:elementary-steps:catalyst}{catalisador}.
\end{proof}

\begin{omfigure}
\setchemfig{atom sep=1.9em}
\schemestart
\chemfig{C(-[:150]R)(-[:210]H)=O}
\arrow{->[\ce{H+}][]}
\chemfig{C^{+}(-[:150]R)(-[:210]H)-OH}
\arrow{->[$\mathrm{R'OH}$][\ce{-H+}]}
\chemfig{C(-[:150]R)(-[:210]H)(-[2]OR')-OH}
\schemestop

\medskip
\schemestart
\chemfig{C(-[:150]R)(-[:210]H)(-[2]OR')-OH}
\arrow{->[\ce{H+}][\ce{-H2O}]}
\chemfig{C^{+}(-[:150]R)(-[:210]H)-OR'}
\arrow{->[$\mathrm{R'OH}$][\ce{-H+}]}
\chemfig{C(-[:150]R)(-[:210]H)(-[2]OR')-OR'}
\schemestop
\par\vspace{6pt}

{\small A formação de um \omterm{def:b1:acetals-protection:hemiacetal}{acetal}, em duas linhas: primeiro o \omterm{def:b1:acetals-protection:hemiacetal}{hemiacetal}
(ativação, adição do álcool, perda de um próton), depois o \omterm{def:b1:acetals-protection:hemiacetal}{acetal}
(protonação do OH, perda de água até o cátion estabilizado, adição de um
segundo álcool, perda de um próton).}
\end{omfigure}

A reação global, $\mathrm{RCHO} + 2\,\mathrm{R'OH} \rightleftharpoons \mathrm{RCH(OR')_2} + \ce{H2O}$, tem
uma constante de equilíbrio modesta, e a água é um produto.

\begin{proposition}[Deslocar a acetalização]\label{prop:b1:acetals-protection:dean-stark}
Se a água formada for removida continuamente, a acetalização prossegue
até que o composto carbonílico (ou o álcool) se esgote.
\end{proposition}

\begin{proof}
O \omterm{def:b1:extent-q-and-k:quotient}{quociente de reação} $Q = [\text{acetal}][\ce{H2O}]/([\text{aldeído}]
[\mathrm{R'OH}]^2)$ permanece abaixo de $K$ enquanto a água for removida:
pelo \cref{ch:b1:extent-q-and-k}, a reação continua avançando, qualquer
que seja o valor de $K$. Com um diol (etano-1,2-diol), o \omterm{def:b1:acetals-protection:hemiacetal}{acetal} é cíclico,
e o equilíbrio também é mais favorável, uma molécula de diol substituindo
duas de álcool.
\end{proof}

\begin{omfigure}
\begin{tikzpicture}[font=\small]
  \pic at (0,0) {heatingmantle};
  \pic at (0,0) {rbflask={0.4}{omDef!14}};
  \pic (d) at (0,3.0) {deanstark={0.35}{omDef!35}{orange!15}};
  \pic at (0,6.4) {condenser};
  \draw[->] (2.4,4.0) node[right, align=left] {coletor: a água (camada\\inferior) não volta mais; o\\tolueno transborda para o balão} -- (0.8,3.8);
  \draw[->] (-2.2,1.0) node[left, align=right] {aldeído, diol, tolueno,\\catalisador ácido, ebulição} -- (-0.7,0.9);
  \draw[->] (1.8,8.0) node[right] {condensador} -- (0.4,7.8);
\end{tikzpicture}

{\small Um aparelho de Dean--Stark. Os vapores de tolueno e de água se
condensam e caem no coletor graduado; a água, mais densa e não \omterm{def:b1:intermolecular-forces-solvents:miscibility}{miscível}
com o tolueno, desce e fica, o tolueno transborda de volta para o balão.
O volume de água recolhido mede o \omterm{def:b1:extent-q-and-k:extent}{avanço da reação}.}
\end{omfigure}

\begin{proposition}[Estabilidade dos acetais]\label{prop:b1:acetals-protection:acetal-stability}
Os \omterm{def:b1:acetals-protection:hemiacetal}{acetais} são estáveis em meio neutro e básico e frente aos \omterm{def:b1:electronic-effects:nucleophile}{nucleófilos},
aos \omterm{def:b1:organomagnesium:organometallic}{reagentes de Grignard} e aos hidretos; o ácido aquoso os hidrolisa de
volta ao composto carbonílico e ao álcool.
\end{proposition}

\begin{proof}
Um carbono acetálico não leva nenhum \omterm{def:b1:electronic-effects:nucleophile}{grupo de saída} que uma base ou um
\omterm{def:b1:electronic-effects:nucleophile}{nucleófilo} possa expulsar: \ce{RO-}, uma \omterm{def:b1:predominant-reaction:strong-acid}{base forte}, não sai, e não resta
nenhuma C=O para atacar. Em ácido aquoso, o mecanismo acima corre ao
contrário: o grande excesso de água torna $Q < K$ para a reação inversa.
\end{proof}

\section{Proteger e desproteger}

\begin{definition}[Grupo protetor]\label{def:b1:acetals-protection:protecting-group}
Um \emph{grupo protetor}\index{grupo protetor} é um grupo introduzido
em uma molécula para mascarar temporariamente uma função, de modo que ela
não reaja em uma etapa dirigida a outra parte da molécula. Sua introdução
é a \emph{proteção}\index{protecao@proteção}, sua remoção a
\emph{desproteção}\index{desprotecao@desproteção}.
\end{definition}

\begin{method}[Proteger, reagir, desproteger]\label{met:b1:acetals-protection:protect}
\begin{enumerate}
  \item Identifique a função que reagiria, ou destruiria o reagente, na
        etapa planejada.
  \item Escolha um \omterm{def:b1:acetals-protection:protecting-group}{grupo protetor} estável nessa etapa e removível em
        condições que o resto da molécula tolera (para um aldeído ou uma
        cetona diante de bases, \omterm{def:b1:electronic-effects:nucleophile}{nucleófilos} ou hidretos: um \omterm{def:b1:acetals-protection:hemiacetal}{acetal} cíclico).
  \item Proteja; realize a etapa planejada; desproteja.
  \item Conte o custo: duas etapas a mais, cada uma com seu \omterm{def:b1:extent-q-and-k:fractional-extent}{rendimento}.
\end{enumerate}
\end{method}

\begin{example}[Uma cetona no caminho]\label{ex:b1:acetals-protection:ketoester}
Para adicionar um \omterm{def:b1:organomagnesium:organometallic}{reagente de Grignard} ao grupo aldeído de uma molécula
que também leva uma cetona, não existe uma escolha simples de \omterm{def:b1:acetals-protection:hemiacetal}{acetal} (os
aldeídos formam \omterm{def:b1:acetals-protection:hemiacetal}{acetais} mais facilmente que as cetonas, de modo que o
grupo errado seria protegido); é preciso mudar a ordem das etapas ou os
reagentes. A \omterm{def:b1:acetals-protection:protecting-group}{proteção} é uma estratégia, não um reflexo.
\end{example}

\section{Exercícios}

\begin{exercise}[$\star$]\label{exo:b1:acetals-protection:1}
Identifique os carbonos de \omterm{def:b1:acetals-protection:hemiacetal}{hemiacetal}, \omterm{def:b1:acetals-protection:hemiacetal}{acetal}, hidrato, éter ou álcool em:
\ce{CH3CH(OH)OCH3}, \ce{CH3CH(OCH3)2}, \ce{CH3CH(OH)2}, \ce{CH3OCH3},
\ce{(CH3)2C(OCH2CH3)2}.
\end{exercise}

\begin{exercise}[$\star$]\label{exo:b1:acetals-protection:2}
Escreva a equação global da formação do \omterm{def:b1:acetals-protection:hemiacetal}{acetal} do propanal com o
metanol, e do \omterm{def:b1:acetals-protection:hemiacetal}{acetal} cíclico da propanona com o etano-1,2-diol.
\end{exercise}

\begin{exercise}[$\star$]\label{exo:b1:acetals-protection:3}
Quais destes reagentes deixam um \omterm{def:b1:acetals-protection:hemiacetal}{acetal} inalterado: solução de hidróxido
de sódio, brometo de metilmagnésio, ácido sulfúrico diluído em água,
boro-hidreto de sódio?
\end{exercise}

\begin{exercise}[$\star$]\label{exo:b1:acetals-protection:4}
Desenhe o \omterm{def:b1:acetals-protection:hemiacetal}{hemiacetal} cíclico formado pelo 4-hidroxibutanal. Qual é o
tamanho do anel?
\end{exercise}

\begin{exercise}[$\star\star$]\label{exo:b1:acetals-protection:5}
Escreva o mecanismo completo da formação, catalisada por ácido, do
\omterm{def:b1:acetals-protection:hemiacetal}{acetal} dimetílico do etanal, com \omterm{def:b1:electronic-effects:curly-arrow}{setas curvas}, e mostre que o ácido é
regenerado.
\end{exercise}

\begin{exercise}[$\star\star$]\label{exo:b1:acetals-protection:6}
Escreva o mecanismo da hidrólise ácida do
2,2-dimetoxipropano. Por que se usa uma grande quantidade de água?
\end{exercise}

\begin{exercise}[$\star\star$]\label{exo:b1:acetals-protection:7}
Uma acetalização de \qty{0.250}{mol} de uma cetona é acompanhada com um
aparelho de Dean--Stark. Que volume de água deve ser recolhido na
conversão completa (densidade de cerca de \qty{1.0}{g/mL})? Depois de uma
hora, \qty{3.2}{mL} foram recolhidos: qual é a conversão?
\end{exercise}

\begin{exercise}[$\star\star$]\label{exo:b1:acetals-protection:8}
Explique por que um \omterm{def:b1:acetals-protection:hemiacetal}{acetal} cíclico com o etano-1,2-diol se forma de modo
mais completo que o \omterm{def:b1:acetals-protection:hemiacetal}{acetal} com duas moléculas de metanol, para o mesmo
composto carbonílico.
\end{exercise}

\begin{exercise}[$\star\star$]\label{exo:b1:acetals-protection:9}
Por que o \omterm{def:b1:elementary-steps:catalyst}{catalisador} ácido precisa ser removido (ou neutralizado) antes
de um composto protegido ser tratado com um \omterm{def:b1:organomagnesium:organometallic}{reagente de Grignard}?
\end{exercise}

\begin{exercise}[$\star\star\star$]\label{exo:b1:acetals-protection:10}
Proponha uma síntese da 5-hidroxi-5-metil-hexan-2-ona,
\ce{(CH3)2C(OH)CH2CH2COCH3}, a partir da 4-clorobutan-2-ona e da propanona,
usando uma \omterm{def:b1:acetals-protection:protecting-group}{proteção}. Dê cada etapa, seus reagentes e o papel de cada um.
\end{exercise}

\begin{exercise}[$\star\star\star$]\label{exo:b1:acetals-protection:11}
Mostre com o \omterm{def:b1:extent-q-and-k:quotient}{quociente de reação} que um grande excesso de metanol também
desloca a acetalização de um aldeído, e compare esse método com a
remoção da água.
\end{exercise}

\begin{exercise}[$\star\star\star$]\label{exo:b1:acetals-protection:12}
Explique por que o \omterm{def:b1:acetals-protection:hemiacetal}{hemiacetal} cíclico da glicose é atacado, em metanol
ácido, dando um \omterm{def:b1:acetals-protection:hemiacetal}{acetal} metílico (um metilglicosídeo), enquanto os outros
grupos OH da glicose não são convertidos em éteres.
\end{exercise}

\section{Problema: um reagente de Grignard que leva um aldeído}

\begin{problem}[{Problema de fim de semana --- por que o 4-bromobutanal não
pode dar um reagente de Grignard, sua proteção como acetal cíclico com um
aparelho de Dean--Stark, a adição de Grignard à propanona, a desproteção
e o volume de água recolhido}]
\label{pb:b1:acetals-protection:1}
Uma química precisa do 5-hidroxi-5-metil-hexanal, \ce{(CH3)2C(OH)CH2CH2CH2CHO}, e
planeja prepará-lo a partir do 4-bromobutanal, \ce{BrCH2CH2CH2CHO}, e da propanona.
Ela parte de \qty{15.09}{g} de 4-bromobutanal. Massas molares
(\unit{g/mol}): \ce{C4H7BrO} 150.9, \ce{C2H6O2} 62.0, \ce{C6H11BrO2} 194.9,
\ce{C3H6O} 58.0, \ce{C9H18O3} 174.0, \ce{C7H14O2} 130.0, \ce{H2O} 18.0;
massa específica da água \qty{0.997}{g/mL}.
% ledger: rho:water

\medskip
\noindent\textbf{Parte I --- O problema.}
\begin{enumerate}
  \item Escreva o \omterm{def:b1:organomagnesium:organometallic}{reagente de Grignard} que o 4-bromobutanal daria.
  \item Por que esse reagente não pode existir? Escreva a reação que o
        destrói.
  \item Que função deve ser protegida, e contra o quê?
  \item Por que um \omterm{def:b1:acetals-protection:hemiacetal}{acetal} é adequado?
  \item Por que se prefere um \omterm{def:b1:acetals-protection:hemiacetal}{acetal} cíclico (com o etano-1,2-diol)?
  \item Calcule a quantidade de matéria de 4-bromobutanal.
  \item Calcule a massa de diol para um excesso de \qty{20}{\percent}.
\end{enumerate}

\medskip
\noindent\textbf{Parte II --- A \omterm{def:b1:acetals-protection:protecting-group}{proteção}.}
\begin{enumerate}[resume]
  \item Escreva a equação da acetalização com o etano-1,2-diol.
  \item Qual é o papel do \omterm{def:b1:elementary-steps:catalyst}{catalisador} ácido (ácido
        4-metilbenzenossulfônico)?
  \item Descreva o funcionamento do aparelho de Dean--Stark.
  \item Explique com $Q$ e $K$ por que remover a água desloca a reação.
  \item Por que é o tolueno, e não a água, que volta ao balão?
  \item Como a química sabe que a reação terminou?
  \item Calcule a massa de brometo protegido esperada na conversão completa.
\end{enumerate}

\medskip
\noindent\textbf{Parte III --- A etapa de Grignard.}
\begin{enumerate}[resume]
  \item Escreva a formação do \omterm{def:b1:organomagnesium:organometallic}{reagente de Grignard} a partir do brometo
        protegido.
  \item Escreva sua adição à propanona e o \omterm{def:b1:alcohol-activation:alkoxide}{alcóxido} formado.
  \item Por que o tratamento desta etapa deve ser levemente básico ou
        neutro (por exemplo, cloreto de amônio aquoso) em vez de ácido?
  \item Calcule a massa de álcool protegido para um \omterm{def:b1:extent-q-and-k:fractional-extent}{rendimento} de
        \qty{70}{\percent} nesta etapa.
\end{enumerate}

\medskip
\noindent\textbf{Parte IV --- A \omterm{def:b1:acetals-protection:protecting-group}{desproteção}.}
\begin{enumerate}[resume]
  \item Escreva a \omterm{def:b1:acetals-protection:protecting-group}{desproteção} em ácido aquoso.
  \item Por que o álcool terciário sobrevive a um ácido aquoso brando?
  \item O produto pode fechar um anel de seis átomos adicionando seu OH ao
        aldeído. Desenhe esse \omterm{def:b1:acetals-protection:hemiacetal}{hemiacetal} cíclico.
  \item Calcule a massa de hidroxialdeído para um \omterm{def:b1:extent-q-and-k:fractional-extent}{rendimento} de
        \qty{90}{\percent} na \omterm{def:b1:acetals-protection:protecting-group}{desproteção}.
  \item Que \omterm{def:b1:extent-q-and-k:fractional-extent}{rendimento} global a partir do 4-bromobutanal a química atinge?
  \item Calcule a massa de água formada na etapa de \omterm{def:b1:acetals-protection:protecting-group}{proteção} na
        conversão completa.
  \item Indique o volume de água recolhido no aparelho de Dean--Stark na
        conversão completa.
\end{enumerate}
\end{problem}
